\subsection{The rank-three support and its arithmetic transfer}
\label{subsec:selmer-support-rank-three}

Keep the same $\mathcal O,F,C_0$ and every original
\texttt{Ok}/\texttt{Big} hypothesis as in
Section~\ref{subsec:selmer-support-rank-two}.  In particular
$\ell$ is odd, $d>3C_0+7$, $d_v>5$, and the original nonzero
$\Psi^{-1}\delta$ and $\Psi^{-2}\delta$ ordinary classes are retained.
Use the actual normalized spaces $\mathscr S_2,\mathscr S_3$,
entry algebras $A_i$, coefficient algebras $B_i$, and fixed-weight
ideals $\lambda_i$, including the coefficient uniformizer $\varpi$.
Write
\[
 \operatorname{FinLam}(I):\quad
 \mathfrak m_{B_i}\subseteq\sqrt{I+\lambda_i}.
\]
This is the original inclusion, also when $I=B_i$.
The rank-three weights are $(1,0,-1)$, and $E$ is the actual ideal
of the four odd entries $r_{01},r_{10},r_{12},r_{21}$.
A type-E prime contains $E$ and is off both triangular loci;
a type-I prime is off those loci and outside $E$.
All charts retain their specified finite unramified coefficient
extension and residue embedding.  The abstract arguments allow
equal characteristic; the actual arithmetic \texttt{Setup} below
also retains its condition $\ell\ne0$ in $\mathcal O$.
\cite{support-rankthree-source}

\begin{theorem}[The actual type-E supplier]\label{thm:rankthree-type-e}
Suppose $I_2\subset B_2$ satisfies \texttt{FinLam} and is contained
in the coefficient contraction of every original rank-two stable
prime.  There is an ideal $I_3\subset B_3$ satisfying
\texttt{FinLam} and contained in the coefficient contraction of
every original type-E prime.
\end{theorem}
\begin{proof}
Put $C=B_3/(E\cap B_3)$ and $K=A_3/E$, with $C\hookrightarrow K$.
The case $E=A_3$ has the witness $I_3=B_3$.  Otherwise the actual
representation over $K$ is the corner $\sigma$ on $(0,2)$ plus a
middle character $\chi$.  Its form and marking restrict to the
original rank-two form and marking.

The ordered corner equations require an argument even though the
middle factor is invertible.  Choose the original local separator
$s$, and set $M=\sigma(s)-\psi_1(s)$ and
$P=\sigma(x)-\psi_0(x)$.  Its determinant belongs to $C$ and is
a unit, since its residue is $(\Psi(s)-1)(\Psi(s)^{-1}-1)$.
The full rank-three ordered product gives, for every $x,y$,
\[
 PM(\sigma(y)-\psi_2(y))=0,\qquad
 PM^2=(\psi_2(s)-\psi_1(s))PM.
\]
Cayley--Hamilton in rank two therefore gives
\begin{equation}\label{eq:rankthree-corner-deletion}
 \det(M)P=(\operatorname{tr}M-\psi_2(s)+\psi_1(s))PM.
\end{equation}
Right multiplication by $\sigma(y)-\psi_2(y)$ proves the exact
ordered rank-two equation.  The middle entry of the original product,
with its outer elements both $s$, has two unit outer factors; hence
$\chi|_{G_v}=\psi_{v,1}$.  No nonreduced quotient was discarded.

The corner gives compatible maps $\xi:A_2\to K$ and
$\beta:B_2\to C$.  To construct them integrally, complete $K$ at
its graded vertex.  Its generators have degrees $0,\pm2$,
opposite products lie in $\mathfrak m_C$, and its fixed pieces are
finite $C$-modules.  If $\mathfrak v$ is the vertex ideal, then
\[
 \mathfrak m_C^a\subseteq\mathfrak v^a\cap C
       \subseteq\mathfrak m_C^{\lceil a/2\rceil}.
\]
Consequently $C$ is closed and $K$ injects into its vertex completion.
The original universal normalized lift maps $B_2$ into closed $C$
and each fixed entry piece into the original $K$.
This constructs $\xi,\beta$ with the exact outer characters,
without evaluating completed series.

For $p:B_3\to C$ take $I_3=p^{-1}(C\beta(I_2))$.
At a prime over $I_3+\lambda_3$, the inverse image in $B_2$ is
maximal by \texttt{FinLam}.  In the resulting complete
characteristic-$\ell$ domain $D$, the character $\chi$ is
residually one, anticyclotomic, and trivial on inertia:
raise inertia elements to $|k|-1$ and use the simple root at one.
The original class-group condition supplies a fixed nonzero exponent
$h$ on every finite $\ell$-group quotient of this unramified
anticyclotomic character.  Completeness gives $\chi^h=1$ in $D$;
characteristic-$\ell$ domain arithmetic and the prime-to-$\ell$
simple-root argument give $\chi=1$.  Every balanced monomial is
now a product of a middle entry and a balanced corner monomial,
so all of $B_3$ is residual.  This proves \texttt{FinLam}.
The middle character was killed only after the fixed-weight cut.
Finally, at any type-E prime the corner is off both rank-two
triangular loci, so rank-two \texttt{AllGood}, transported by
$\xi,\beta$, kills $I_3$ there.
\end{proof}

\begin{theorem}[The actual determinant section]\label{thm:rankthree-sharp}
Let $R=L.R$ be an original chart at a closed point $x$ of the
coefficient-fibre open $\mathcal B$ where a lower entry is nonzero.
There is an ideal $Q^\sharp\subset R$, zero if $x\notin V(E)$,
contained in every prime outside $E$, such that its minimal primes
have dimension at least $3d+2$, are connected through intersections
of dimension at least $3d+1$, and are either type I or original
\texttt{Thick} type E.
\end{theorem}
\begin{proof}
The original chart represents the deformation functor of its actual
lower-triangular residual representation.  The chart proof used in
rank two applies with the finite Laurent algebra
$C[r(g)_{ij}z^{w_i-w_j}]\subset C[z,z^{-1}]$.
Balanced atoms have length at most three, and upper coefficients
lie in $\mathfrak m_C$, so
\[
 \mathfrak m_C^a\subseteq\mathfrak v^a\cap C
       \subseteq\mathfrak m_C^{\lceil a/3\rceil}.
\]
Fixed pieces are finite complete $C$-modules.  The first nonzero
coefficient modulo powers of $\mathfrak m_C$ proves injectivity
through vertex localization and completion.  Thus the original
universal property maps the coefficient algebra into closed $C$;
evaluation at $z=1$ is used only in the finite Laurent algebra.

Form normalization is unobstructed because $2$ is a unit.  The
matrix $D=r(g_0)-J_3r(g_0)^tJ_3$ is skew-adjoint, with trace and
determinant zero and residual eigenvalues $\zeta_0,0,-\zeta_0$.
Hensel's lemma and the three spectral lines give a strict orthogonal
change of basis making $D=\operatorname{diag}(\zeta,0,-\zeta)$.
Only $2$ and $\zeta$ are inverted.  The remaining gauge is
$\operatorname{diag}(u,1,u^{-1})$.  Its adjoint stabilizer is zero:
the mark makes an invariant diagonal, the nonzero lower entry and
polarization identify the relevant diagonal coordinates, and the
$c$-action $X\mapsto-J_3X^tJ_3$ kills them.  This includes $\ell=3$.

Suppose first $x\in V(E)$.  Its residue is
\[
 r_0(g)=\begin{pmatrix}\Psi(g)&0&0\\0&1&0\\a(g)&0&\Psi(g)^{-1}\end{pmatrix},
 \qquad a\ne0.
\]
Conjugation by $\operatorname{diag}(1,-1,1)$ acts on the actual
chart and splits the full unitary adjoint into dimensions $5$ and $4$.
The unique local ordered flags are invariant under this involution.
Indeed, at a separating local element the three spectral projectors
$E_i$ have rank one.  Leave only factor $k$ variable in the complete
ordered product and multiply by $E_i,E_j$: the remaining differences
are units for $i\ge k\ge j$.  This gives the required diagonal
characters and kills every entry below the prescribed flag, over
the nonreduced Artinian ring as well.

The local even and odd Borels have dimensions $4,2$, and the real
adjoint invariants have dimensions $1,2$.  The even constituents are
$\delta,\Psi^{-2}\delta,1,\Psi^2\delta,\delta$; the odd constituents
are $\Psi^{-1}\delta,\Psi^{-1},\Psi\delta,\Psi$.
The original \texttt{f2a}/\texttt{f2c} and \texttt{f2b} respectively
exclude every cyclotomic quotient.  Thus the exact Euler inequality
for the original continuous ordinary cone gives
\begin{equation}\label{eq:rankthree-parity-euler}
 h^1_e-h^2_e\ge3d,\qquad h^1_o-h^2_o\ge0.
\end{equation}
The global $H^0$ groups vanish as above.  The orthogonal gauge has
one even and two odd directions; the two odd marking equations
remove the latter by multiplication by units $\pm\zeta_0$.
The normalized tangent dimensions are therefore $a=h^1_e+1$ and
$g=h^1_o$.  The same multiplication-and-local-flag obstruction
calculation as in rank two gives $H^2_e,H^2_o$; form and marking
normalization add no obstruction.  Since parity projection is exact,
the dual minimal relation spaces inject separately into these two
groups.  Choose parity variables and relations and pad with zero
relations.  We obtain the actual presentation
\begin{equation}\label{eq:rankthree-parity-presentation}
 R=\mathcal O'[[X_1,\ldots,X_a,Y_1,\ldots,Y_g]]/(F_1,\ldots,F_b,G_1,\ldots,G_g),
 \quad a-b=3d+1,
\end{equation}
where $F$ is even, $G$ is odd and every linear relation coefficient
lies in $\mathfrak m_{\mathcal O'}$.  The actual ideal $E$ is $(Y)$:
the involution is identity modulo $E$, while every odd series is
divisible by an odd variable.  Nilpotents are retained.

Here $g>0$ follows from the original arithmetic nonvanishing.
The lower odd submodule $U=\langle E_{10},E_{21}\rangle$ is
$k'(\Psi^{-1}\delta)\oplus k'(\Psi^{-1})$; its quotient has
$G_0$-character $\Psi$ and no invariants.  If the first local
corner line is $e_0+\lambda e_2$, the odd flag stabilizer is
$X_{10}+\lambda X_{12}=0$, $X_{21}=\lambda X_{01}$.
It meets $U$ in zero and maps isomorphically onto the quotient.
The exact sequence of the original ordinary cones therefore injects
\[
 H^1_C(\Gamma,k'(\Psi^{-1}\delta);0)\hookrightarrow H^1_o.
\]
The original \texttt{f4} makes the source nonzero.

Write $G=MY$ with $M$ even and constant coefficients in
$\mathfrak m_{\mathcal O'}$, and put $\Delta=\det M$.
The adjugate identity gives $\Delta Y_j=0$ in $R$.
For connectedness let $A_0=\mathcal O'[[X,Y]]$ and consider the
generic incidence ideal
\begin{equation}\label{eq:rankthree-incidence}
                  (Zv,\det Z)\subset A_0[Z_{ij},v_j].
\end{equation}
On $v_i\ne0$ its equations solve one column of $Z$, giving a
polynomial domain; these charts overlap.  Every point with $v=0$
and singular $Z$ is in their closure, by the family $(Z,tu)$ for
a nonzero kernel vector $u$ over its residue field.  Thus the
radical is prime, with relative dimension $g^2$.
It stays prime after vertex completion: the homogeneous quotient
has finite complete $A_0$-pieces, and the first nonzero homogeneous
terms of two series have nonzero product.  A power of the radical
lies in the original ideal, so both completions have the same
prime radical.  The complete incidence ring has dimension
$a+g+1+g^2$.  Cutting by $Z-M$, $v-Y$ and $F$ uses $g^2+g+b$
elements and gives exactly $R/(\Delta)$.  Hence all its components
have dimension at least $3d+2$.  Complete-local connectedness
\cite{support-ranktwo-connectedness}, with cohomological dimension
bounded by this number of cuts, connects them after deletion of
any closed set of dimension $<3d+1$.  This is the stated intersection
graph, without a regular-sequence or transversality hypothesis.

If $P\supset(E,\Delta)$ is prime, put $Q=R/P$ and choose
$0\ne\nu\in\ker(M:Q^g\to Q^g)$ by clearing finitely many
fraction-field denominators.  The derivation
\[
 D_0=\sum_j\nu_j\frac{\partial}{\partial Y_j}\pmod P
\]
kills $F$ and $MY$, hence descends to $D:R\to Q$.  It kills all
even entries and characters, but some odd entry has nonzero image.
The actual $Q$-valued representation
\begin{equation}\label{eq:rankthree-thick-matrix}
 \widehat r(g)=
 \begin{pmatrix}\bar r_{00}(g)&D r_{01}(g)&\bar r_{02}(g)\\
 0&\bar r_{11}(g)&0\\\bar r_{20}(g)&D r_{21}(g)&\bar r_{22}(g)\end{pmatrix}
\end{equation}
is multiplicative by Leibniz.  The same block operation applied to
the original ordinary product, and to its polarized star reversal,
gives both original ordered products.  No polarization of
$\widehat r$ is assumed.  A splitting over any injective field
extension would have eigenvector $(z_0,1,z_2)$.  Subtracting its
equations at $g_0,g_0^*$ gives $\zeta z_0=\zeta z_2=0$;
the eigenvector equations then kill its odd column, and the original
polarization kills the other odd derivatives, a contradiction.
Finally $D(\mathfrak m_R^{N+1})\subset\mathfrak m_Q^N$ proves
the required adic continuity.  This is literal \texttt{Thick}.

For the needed dimension bound, at a prime of the coefficient fibre
the invertible lower entry forces every upper entry to vanish:
an upper of degree $b$ and lower of degree $-a$ have balanced
product $\mathrm{upper}^a\mathrm{lower}^b\in\mathfrak m_B$.
The matrix is consequently
\[
 \begin{pmatrix}\Psi&0&0\\y&1&0\\a&y\circ *&\Psi^{-1}\end{pmatrix},
\]
with the exact equations
\begin{align}
 y(gh)&=\Psi(h)y(g)+y(h),\notag\\
 a(gh)&=\Psi(h)a(g)+y(g^*)y(h)+\Psi(g)^{-1}a(h),\notag\\
 a(g^*)&=a(g),\quad y(g_0)=y(g_0^*),\notag\\
 y_v&=(\Psi-1)\alpha_v,\quad y_v^*=(1-\Psi^{-1})\beta_v,\notag\\
 a_v&=(\Psi-\Psi^{-1})\gamma_v-(1-\Psi^{-1})\alpha_v\beta_v.
 \label{eq:rankthree-low}
\end{align}
The original \texttt{b1} bounds the $y$-space by $C_0$:
$y/\Psi$ is its cocycle, local trivializations are unique, and
a kernel coboundary $t(\Psi-1)$ is killed by the marking.
For fixed $y$, differences of $a$ lie in the strict
$\Psi^{-2}\delta$ space, bounded by $C_0$ by \texttt{b1p}.
Entry functions in the coefficient fibre have finite image, by the
finite fixed-piece topology.  At a generic point, over the field
generated by the finitely many $y$-values the equations on the
$a$-values are linear.  Solvability after extension implies
solvability over that field, and their homogeneous solution space
has dimension at most $C_0$.  The finite-type closed-point dimension
formula and the specified completion give
\begin{equation}\label{eq:rankthree-bad-dimensions}
 \dim R/\mathfrak m_BR\le2C_0,
 \qquad \dim R/T_-\le2d+3C_0+2<3d.
\end{equation}
For the second inequality cut $R/T_-$ by $\varpi$, at most $C_0+1$
norm-character detectors, and $d$ inertia generators for each of
the first and middle diagonal characters.  The last character is
determined by polarization.  The original generator/class-group
conditions make all diagonal and ordered local characters residual,
so these $2d+C_0+2$ cuts kill $\mathfrak m_B$ up to radical.
The height theorem gives the bound.  This argument uses prime
quotients for dimension; upper entries in the actual fibre can
still be nilpotent.

The lower triangular ideal is a unit in this chart, and
\eqref{eq:rankthree-bad-dimensions} excludes the other triangular
locus from every minimal component of $R/(\Delta)$.
Those outside $E$ are type I and those containing $E$ are Thick.
Thus $Q^\sharp=(\Delta)$ works.  If $x\notin V(E)$, use
$Q^\sharp=0$: the full adjoint has local Borel dimension $6$,
real invariant dimension $3$, and one remaining gauge direction,
so its tangent-minus-relation bound is $3d+1$.  The same height
and connectedness argument gives dimensions $3d+2$ and $3d+1$;
$E=R$ and the triangular bound make every component type I.
\end{proof}

\begin{theorem}[Complete normalized rank-three propagation]
\label{thm:rankthree-normalized-propagation}
Assume the rank-two input of Theorem~\ref{thm:rankthree-type-e}.
For $I\subset B_3$ retain the original \texttt{FinLam}, P1', P23'
and weak starting-support condition
\begin{equation}\label{eq:rankthree-weak-seed}
 I\subset Z\subset\xi\cap B_3,\qquad \xi\text{ type I},
                    \dim(B_3/Z)=3d+1.
\end{equation}
Then $I$ is contained in every type-I coefficient prime, and
$I\cap I_3$ is a \texttt{FinLam} ideal contained in every
coefficient prime off the two triangular loci.
\end{theorem}
\begin{proof}
We first supply the literal Nice curves.  Let $H$ satisfy
\texttt{FinLam} and $P\supset HR$ have $\dim R/P\ge3d+1$.
For a prime off the triangular loci, nonzero entries of degrees
$a,-b$ give a nonzero balanced product $u^bv^a$.  Every entry $z$
of degree $w$ satisfies $z^au^{-w}\in\operatorname{Frac}(B/\mathfrak p)$,
after clearing powers of that balanced product.  Thus the algebraized
entry field has relative transcendence degree at most one.
The actual closed-point dimension formula over complete $B/\mathfrak p$,
followed by the prescribed finite etale extension, completion and
prime quotient, gives $\dim R/Q\le\dim B/\mathfrak p+1$.
This makes no assertion about the transcendence degree of a completed
fraction field.

At a coincidence $\psi_{v,i}|_{N_v}=\psi_{v,j}|_{N_v}$, the original
inertia coordinates save $d_v-1$ generators among the three sets
of $d$ coordinates.  Including $\varpi$, at most $3d-d_v+2$
cuts kill the fixed-weight fibre; \texttt{FinLam} and the preceding
bound give
\[
                         \dim R/Q\le3d-d_v+3<3d.
\]
Together with \eqref{eq:rankthree-bad-dimensions}, this allows the
rank-two prime-avoidance argument: cut $P$ by $\varpi$, then to
dimension one while avoiding all triangular and ordered coincidence
loci.  The result is the original characteristic-$\ell$ Nice curve.
If it contains $E$ it is type E; otherwise its marked projectors
give absolute irreducibility.  Indeed every invariant coordinate
line or plane either forces triangularity or is the middle line
or corner plane, which forces $E$.

The original open coefficient fibre $\mathcal B$ is connected.
Its field points are the nonzero pairs in
\eqref{eq:rankthree-low}.  These polynomial families map to the
original $A_3$: introduce finitely many entry values and local
parameters through a common finite quotient of $G_0$, impose the
weighted homogeneous Low equations, and use their vertex completion
in the original universal property.  Each homogeneous entry maps
to the uncompleted polynomial algebra, and $B_3$ maps residually.
Evaluation is only in this polynomial algebra.
Choose a nonzero strict corner cocycle $b$ from the original
\texttt{f4'}.  The Low family
$(Ty,T^2a+(1-T)b)$ connects the generic point to both $(y,a)$
and $(0,b)$, so every point is connected to the latter.

Moreover $\mathcal B\setminus V(E)$ is dense.  From \texttt{f4}
choose a nonzero $\Psi^{-1}\delta$ cocycle and write its normalized
$G_0$ version as $u$, with
$u(gh)=\Psi(h)u(g)+u(h)$ and $u(g^*)=\Psi(g)^{-1}u(g)$.
Choose $c$ by $c(\Psi(g_0)^2-1)=-u(g_0)(\Psi(g_0)-1)$ and set
\begin{align*}
 y_0&=u+c(\Psi-1),\\
 a_0&=c^2-\tfrac12c^2(\Psi+\Psi^{-1})
             -c(1+\Psi^{-1})u+\tfrac12\Psi^{-1}u^2.
\end{align*}
Direct substitution gives every Low equation; the local parameters
are $\alpha+c,\alpha-c,\alpha^2/2-c\alpha-c^2/2$.
The marking and $u\ne0$ imply $y_0\ne0$.  At any E point $(0,a)$,
the family $(Ty_0,T^2a_0+a)$ specializes from outside E to that point.

Apply the component graph of Theorem~\ref{thm:rankthree-sharp}.
For an edge between type-I components, a Nice curve above their
intersection transfers goodness by P1', or by the length-zero P23'
case if the curve is type E.  For a consecutive block of Thick
type-E components, choose Nice curves above consecutive intersections
using the separate finite-weight ideal $I_3$.  Consecutive curves
contain their common actual Thick component.  The full P23' chain
therefore transfers goodness from the type-I entrance to the type-I
exit; intermediate type-E components need not contain $I$.

The weak seed \eqref{eq:rankthree-weak-seed} supplies the first chart.
Choose at $\xi$ a nonzero odd entry $u$ of degree $\pm1$ and a
nonzero opposite entry $v$ of degree $\mp b$, $b\in\{1,2\}$.
For $a=u^bv\notin\xi$ the actual grading gives
\[
                     (A_3)_a=(B_3)_a[u,u^{-1}].
\]
Thus $Z(B_3)_a[u,u^{-1}]$ contracts to a homogeneous type-I prime
with coefficient prime exactly $Z$.  It meets $\mathcal B$:
otherwise a power of its lower-entry ideal would lie in
$\mathfrak m_B$, and Nakayama on sufficiently negative finite
graded pieces would kill powers of a nonzero lower entry.
At a closed point there the original chart has a prime of dimension
$\dim(B_3/Z)+1=3d+2$ with this contraction.  Select a maximal-dimension
minimal prime of the completion; flatness preserves that contraction.
The Nice-curve argument gives a good minimal type-I component there.

Every type-I component meets $\mathcal B$ by the same graded argument.
At a closed point of $\mathcal B\setminus V(E)$ the full-adjoint
presentation bounds every local component below by $3d+2$, while
the triangular locus has dimension less than $3d$.  Thus every
local component is type I.  Faithful flatness and the density just proved show that
the union of type-I components covers $\mathcal B$.
Local propagation makes the intersections with the finite unions
of good and bad type-I components disjoint closed subsets: a
nonempty intersection would contain a closed point of this finite-type
fibre.  Connectedness and the starting component make all components
good.  Every type-I prime contains one of them, proving \texttt{AllGood}.
Finally a prime over $(I\cap I_3)+\lambda_3$ contains $I$ or $I_3$
and is maximal.  Every off-triangular prime is type I or E, so it
contains $I\cap I_3$.  This is exactly \texttt{OffFin}.
\end{proof}

\begin{theorem}[Transfer to the actual arithmetic support]
\label{thm:rankthree-arithmetic-transfer}
Retain the complete original arithmetic \texttt{Setup} and
\texttt{ModIn} inputs in ranks two and three.  The original
determinant map $\Phi:B_{\rm det}\to B_3$ is surjective, its Hecke
ideal image satisfies \texttt{FinLam}, and the original Nice/Thick
transfers supply P1' and the entire P23' chain condition.
Consequently Theorem~\ref{thm:rankthree-normalized-propagation}
applies to the actual rank-three Hecke ideal.
\end{theorem}
\begin{proof}
The same twisted determinant and ordered characters give $\Phi$
by their original universal property.  Its image is closed by
compactness.  Undoing $\kappa$ along the same quotient $\pi$ puts
all traces and ordered characters in this image.  The mark gives
\[
 \zeta^2=\operatorname{tr}(r(g_0)^2-r(g_0)r(g_0^*)),\quad
 e_0=\frac{D(D+\zeta)}{2\zeta^2},\quad
 e_1=1-\frac{D^2}{\zeta^2},\quad
 e_2=\frac{D(D-\zeta)}{2\zeta^2}.
\]
Hensel gives $\zeta$ in the image.  The balanced atoms have types
$(1,-1)$, $(2,-2)$, $(1,1,-2)$ and $(-1,-1,2)$; polarization
$r_{01}(g)=r_{12}(g^*)$, $r_{10}(g)=r_{21}(g^*)$ turns each into
a closed matrix cycle, hence a trace with these projectors.
Thus the image contains all of $B_3$, without a denominator $3$.
The original finite weight-ring condition now gives
\texttt{FinLam} exactly as in rank two: after the uniformizer cut
the inertia characters are their prescribed residual characters,
and a finite coefficient quotient over the residual weight field
is a field.

Normalize any original Nice curve.  Excellence and completeness
give a finite complete DVR $k'[[t]]$ with finite residue field,
the same kernel and a local map from its chart.  The original
$\kappa$ twist gives the original Point: continuity and residual
determinant are preserved, $J_3$ supplies its symmetric form,
and the spectral-projector ordered-product argument gives its
local flags.  The inertia verification is included below.
On each original $N_v$ the ratios of local characters are
residually one and nontrivial, so have infinite order in
$1+tk'[[t]]$.  They remain nontrivial in any prescribed open subgroup.
Choose finitely many witnesses; their image in the product of
ratio groups is finitely generated torsion-free abelian.
Avoiding its finitely many proper rational kernel hyperplanes gives
one local element $h\in U\cap G'$ with all eigenvalue ratios of
infinite order.  The exact TwLink identity gives $\epsilon(h)=1$.
This proves both original N1 inequalities simultaneously.

For type I, the representation stays absolutely irreducible on
every $U\cap G'$.  Replace $U$ by its normal core.  The radical
of the finite-dimensional restricted image algebra is stable under
the full irreducible group and nilpotent, hence zero.  Its isotypic
summands are permuted transitively.  In dimension three a reducible
restriction is scalar, excluded by $h$, or three distinct lines
permuted transitively.  In the latter case each local group fixes
all three lines: it preserves a full ordered flag and contains
the separating $h$ in the normal subgroup.  Let $H$ be the permutation
kernel and $g$ a 3-cycle.  Every $r(gh)$ for $h\in H$ has trace zero,
so residual reduction and cancellation of $\bar\kappa$ give
\[
                    \Psi(gh)+1+\Psi(gh)^{-1}=0.
\]
In characteristic three this forces $\Psi(g)=\Psi(H)=1$.
Otherwise $\Psi(g)=\zeta_3\ne1$ and
$\Psi(H)\subset\{1,\zeta_3^{-2}\}$; this set contains no nontrivial
subgroup, so again $\Psi(H)=1$.  Since local groups lie in $H$,
both cases contradict the original local condition \texttt{f3}.
For type E the actual corner remains irreducible on every $U\cap G'$:
the corresponding two-line alternative gives
$\Psi(gh)^2=-1$ outside its permutation kernel, hence $\Psi(h)^2=1$
on that kernel and on every local group, again contrary to \texttt{f3}.

Here is the full N3 argument, including type E and $\ell=3$.
For $G_1=V\cap G'$ let $H$ be the reduced Zariski closure of its
image over an algebraic closure $\Omega$ of $k'((t))$.
The inverse image of $H^0$ is closed of finite index, hence open.
Closedness follows by expanding the finitely many coefficients
in its equations in a basis over $k'((t))$; this reduces to
continuous polynomial equations over that field.
The preceding all-open irreducibility makes $H^0$ irreducible on
the full space, or on the corner in type E.  It is an irreducible
variety.\cite{support-rankthree-torus}
Diagonalize the separating $u\in G_1$.  For any prescribed coset
matrix $Y$, choose $g\in G_1$ such that all diagonal entries of
$M=Yr(g)$ and $M^{-1}$ are nonzero: each is a nonzero
matrix-coefficient function on $H^0$, and their open loci meet.
In type E use its irreducible corner; the middle entry is already
a unit.  For rank at most three these entries and $\det M$ give
all principal minors.

The identity component of the closure of $\langle r(u)\rangle$
is a torus on which every ratio character is nontrivial.  Choose
a cocharacter $\operatorname{diag}(z^{w_1},\ldots,z^{w_n})$ with
distinct integer weights by avoiding finitely many lattice
hyperplanes.\cite{support-rankthree-torus}  Order those weights.
The coefficient of $X^{n-j}$ in its product with $M$ has valuation
$w_1+\cdots+w_j$, since its unique lowest-weight principal minor
is nonzero.  Substituting $X=z^{w_j}Z$ and dividing by the least
power of $z$ gives reduction
$Z^{n-j}(a_{j-1}Z+a_j)$ with $a_{j-1}a_j\ne0$.
Its nonzero root is simple in every characteristic; Hensel gives
a root of valuation $w_j$.  Hence all $n$ roots are distinct.
The discriminant is nonzero on that torus, so some power $u^m$
gives separable $Yr(gu^m)$.  For an arbitrary open set meeting $G'$,
choose $s\in G'$ in it and $V$ with $sV$ inside it; taking $Y=r(s)$
proves precisely N3.

Type E has the actual $(0,2;1)$ block decomposition.  The original
polarization gives both block forms with multiplier $\epsilon^{-2}$,
the exponent prescribed by rank three.  Its corner is irreducible
on $G'$, its line is irreducible, and every intertwiner from the
corner to the line is zero by the invariant-kernel argument.
These are all original TypeE fields.  Thus the literal Nice transfer
is proved.

For a Thick prime $P$, put $Q=R/P$ and twist the actual modification
\eqref{eq:rankthree-thick-matrix}, obtaining $\sigma$.
Its determinant law is the original twisted determinant law:
over every scalar-extended group algebra its matrices have the same
diagonal corner and middle blocks.  Thus the full multivariable
Cayley--Hamilton coefficient relations vanish.  The ordinary
products vanish by the original ordered equations.  The identity
\[
 \sigma(\operatorname{star}_{\rm old}g)
       =\kappa(g)\widehat r(\operatorname{star}_{\rm new}\pi(g))
\]
follows from the exact TwLink polarization and gives the second
product with reversed order.  For original distinguished inertia,
TwLink gives $m$ prime to $\ell$ with $\pi(g)^m=\kappa(g)^m=1$.
The matrix $M$ for $g$ or its star has $M^m=1$ and residual
characteristic polynomial $(X-1)^3$.  The geometric sum
$1+M+\cdots+M^{m-1}$ is a unit, being residually $m$ plus a
commuting nilpotent.  Thus $M=1$.  This also verifies the Point
inertia condition above.

The actual coefficient map to $Q$ is local: a source nonunit with
nonzero residue in a finite field of size $q$ would make the source
unit $1-a^{q-1}$ map to zero.  Its locality and the adic continuity
of $\sigma$ kill the original closed group-algebra ideal: reduce
modulo every maximal-ideal power, choose an open normal subgroup
where $\sigma$ is trivial, and use Krull intersection in $Q$.
Embed the prescribed algebraically closed coefficient residue field
$\kappa_P$ into an algebraically closed field $L$ containing $Q$
injectively.  The resulting representation factors through the
actual algebra $E_0(\kappa_P)$ and has invariant corner plane
$W=\{v:v_1=0\}$.

The coefficient prime is not stable.  If an irreducible representation
over $\kappa_P$ had the same determinant, Burnside's theorem would
give its full image algebra $M_3(\kappa_P)$.  The middle character
of $\sigma$ factors through this image: an element in its kernel
has characteristic polynomial $X^3$ under $\sigma$, so its middle
entry is zero.  But matrix units exclude a unital homomorphism
$M_3(\kappa_P)\to L$.
Nor is $E_0(\kappa_P)$ semisimple.  Otherwise $W$ has a module
complement containing $v$ with $v_1=1$.  Each middle eigenvalue
lies in $\kappa_P$, because it is a root of its universal
characteristic polynomial over that algebraically closed field.
Thus $\sigma(g)v-\sigma(g)_{11}v$ lies both in the complement and
in $W$, and is zero.  Canceling $\kappa$ and using surjectivity
of $\pi$ gives the forbidden splitting of $\widehat r$ over $L$.
The complement need not be $L$-linear; the coefficient-field
eigenvalue argument is what supplies scalar stability.  These
are exactly the two original \texttt{IsThick} clauses.

A type-I prime contracts to a stable coefficient prime: its twisted
irreducible realization over an algebraically closed extension
descends in the determinant sense to the prescribed algebraic closure
of the residue field, by realization and uniqueness of semisimplification
over algebraically closed fields.\cite{support-chenevier}
Finally apply the unchanged arithmetic P1core to the actual NiceI
curve, using \texttt{ModIn.n4}.  For P23', normalize every member
of the complete Nice type-E chain.  They are NiceE by the same N4
input, their common Thick primes contract to \texttt{IsThick}
primes, and their initial and final type-I primes contract to stable
ones.  These are precisely the hypotheses and terminal alternative
of P3core, including chain length zero.  Contraction along the
same $\Phi$ gives P1' and P23'.
The already proved seed transfer supplies
\eqref{eq:rankthree-weak-seed}, completing the application.
\end{proof}

\paragraph{Scope.}
These theorems discharge the original normalized rank-three component
construction and its arithmetic interface under the actual
\texttt{ModIn} suppliers.  They do not supply a uniform nilpotence
exponent or the moving-family integral Selmer estimate, and do not
by themselves complete the RH or GRH proof-map node.
